\chapter{Support Vector Machines and the Kernel Trick}
\companion{StatQuest: Support Vector Machines Part 1, Main Ideas}{\ytid{efR1C6CvhmE}}

InfoEdge interviews reportedly asked ``what is an SVM?'' and expected, in order: hinge loss, kernels, margin and support
vectors. The OA asked SVM statements and that hinge loss is a classification loss. The ``kernel methods'' line on the
Tools slide shows the team uses them. This chapter builds SVMs slowly enough that you can derive the margin and
explain the kernel trick to someone who has never seen it.

\section{Which line is best?}
Two classes of points can be separated by infinitely many lines. Which should we choose? Intuition: the line that stays as
far as possible from both classes. New points are slightly different from training points; a line hugging one class would
misclassify them at the first wobble. The distance from the line to the nearest points is the \textbf{margin}. An SVM picks
the line with the \textbf{maximum margin}: the widest empty ``street'' between the classes.

\begin{center}
\begin{tikzpicture}[scale=0.9]
\draw[->] (0,0)--(6.5,0); \draw[->] (0,0)--(0,4.5);
\draw[very thick,accent] (0.6,4.2)--(5.85,0);
\draw[dashed,accent] (0.0,3.58)--(4.475,0);
\draw[dashed,accent] (1.5,4.58)--(6.4,0.66);
\foreach \p in {(0.7,1.2),(1.3,0.6),(2.0,1.0),(1.0,2.78),(2.6,1.5)} \fill[prframe] \p circle (2.3pt);
\foreach \p in {(4.2,3.0),(5.0,2.2),(4.8,3.6),(3.8,2.74),(5.8,3.0)} \fill[exframe] \p circle (2.3pt);
\draw[thick] (2.6,1.5) circle (4.5pt); \draw[thick] (1.0,2.78) circle (4.5pt); \draw[thick] (3.8,2.74) circle (4.5pt);
\node[font=\small,align=left] at (8.6,2.2) {solid: $\mathbf w^\top\mathbf x + b = 0$\\dashed: $\mathbf w^\top\mathbf x + b = \pm1$\\circled: support vectors\\street width $= 2/\norm{\mathbf w}$};
\end{tikzpicture}
\end{center}

Only the points on the edges of the street matter. They are the \textbf{support vectors}. Move or delete any other point
(without crossing the street) and the solution does not change.

\section{The geometry you need}
A hyperplane is $\mathbf w^\top\mathbf x + b = 0$. The vector $\mathbf w$ is perpendicular to it. The signed distance from a point
$\mathbf x$ to the plane is
\[ \frac{\mathbf w^\top\mathbf x + b}{\norm{\mathbf w}}. \]
\begin{example}[: distance to a plane]
$\mathbf w = (3, 4)$, $b = -5$, point $(1, 1)$: $\frac{3 + 4 - 5}{\sqrt{9+16}} = \frac{2}{5} = 0.4$.
\end{example}
Labels are $y\in\{-1,+1\}$. A point is correctly classified when $y(\mathbf w^\top\mathbf x + b) > 0$.

\section{Hard margin: the optimisation problem}
Scaling $\mathbf w$ and $b$ by the same constant does not move the plane. Use that freedom: scale so the closest points satisfy
$y(\mathbf w^\top\mathbf x + b) = 1$ exactly. Then their distance to the plane is $1/\norm{\mathbf w}$ and the street width is
$2/\norm{\mathbf w}$. Maximising $2/\norm{\mathbf w}$ is the same as minimising $\frac12\norm{\mathbf w}^2$ (squared for smoothness):
\[ \min_{\mathbf w,b}\ \tfrac12\norm{\mathbf w}^2\quad\text{subject to}\quad y_i(\mathbf w^\top\mathbf x_i + b)\ge1\ \text{for all } i. \]
This is a convex quadratic program: a bowl-shaped objective with linear constraints, so there is one global solution.

\begin{example}[: the simplest SVM]
Two points: $(1,1)$ with $y = +1$ and $(-1,-1)$ with $y = -1$. By symmetry the boundary is perpendicular to the segment joining
them and passes through the origin: $\mathbf w = (a, a)$, $b = 0$. Constraint for $(1,1)$: $2a \ge 1$, smallest $\norm{\mathbf w}$ at
$a = 0.5$. So $\mathbf w = (0.5, 0.5)$, margin $2/\norm{\mathbf w} = 2/0.707 = 2.83$, which is exactly the distance between the two
points. Both points are support vectors.
\end{example}

\section{Soft margin: allowing mistakes, and the C parameter}
Real data overlap, and one outlier could force a tiny margin. Introduce a \textbf{slack} $\xi_i\ge0$ per point that measures how
much it violates the margin:
\[ \min\ \tfrac12\norm{\mathbf w}^2 + C\sum_i\xi_i\quad\text{s.t.}\quad y_i(\mathbf w^\top\mathbf x_i + b)\ge1 - \xi_i,\ \xi_i\ge0. \]
\begin{itemize}
\item $\xi_i = 0$: outside the street on the correct side. $0<\xi_i\le1$: inside the street, still correct.
  $\xi_i > 1$: misclassified.
\item \textbf{Large $C$}: violations are expensive; narrow margin; fits training data hard; low bias, high variance.
  $C\to\infty$ is the hard margin.
\item \textbf{Small $C$}: wide margin; tolerates violations; more bias, less variance.
\end{itemize}

\subsection{The same thing as a loss function: hinge loss}
At the optimum, $\xi_i = \max(0, 1 - y_if(\mathbf x_i))$. Substituting gives an unconstrained problem:
\[ \min_{\mathbf w,b}\ \tfrac12\norm{\mathbf w}^2 + C\sum_i\max\big(0, 1 - y_if(\mathbf x_i)\big). \]
The term $\max(0, 1 - yf)$ is the \textbf{hinge loss}. So a linear SVM is ``hinge loss + L2 regularisation'', just as logistic
regression is ``log loss + L2''.
\begin{example}[: hinge loss values]
$y = +1$: $f = 2$ gives 0 (confidently right, ignored); $f = 0.3$ gives $0.7$ (right but inside the margin); $f = -2$ gives 3 (wrong).
$y = -1$, $f = -1.5$: $yf = 1.5$, loss 0.
\end{example}
Hinge loss is exactly zero beyond the margin: confidently correct points exert no pull. That is why only support vectors matter.

\section{The dual problem: where kernels come from}
Using Lagrange multipliers $\alpha_i\ge0$ (one per constraint), the solution has the form
\[ \mathbf w = \sum_i\alpha_iy_i\mathbf x_i, \]
and the $\alpha$'s solve the \textbf{dual} problem
\[ \max_{\boldsymbol\alpha}\ \sum_i\alpha_i - \tfrac12\sum_{i,j}\alpha_i\alpha_jy_iy_j\,\mathbf x_i^\top\mathbf x_j\quad\text{s.t.}\quad 0\le\alpha_i\le C,\ \sum_i\alpha_iy_i = 0. \]
Three facts follow (from the KKT conditions):
\begin{itemize}
\item $\alpha_i = 0$ for points strictly outside the margin: they do not appear in $\mathbf w$.
\item $0 < \alpha_i < C$: points exactly on the margin.
\item $\alpha_i = C$: points inside the margin or misclassified.
\end{itemize}
And the crucial observation: \textbf{the data appear only through dot products} $\mathbf x_i^\top\mathbf x_j$, both in training and in
prediction, $f(\mathbf x) = \sum_i\alpha_iy_i\,\mathbf x_i^\top\mathbf x + b$.

\section{The kernel trick}
\subsection{The problem: curved boundaries}
Points inside a circle (class $+1$) and outside it (class $-1$) cannot be separated by a line. But map each point
$\mathbf x = (x_1, x_2)$ to a new 3-D feature vector $\phi(\mathbf x) = (x_1^2, \sqrt2x_1x_2, x_2^2)$. The circle $x_1^2 + x_2^2 = r^2$
becomes the \emph{plane} $z_1 + z_3 = r^2$. Linear separation in the new space is a curved boundary in the old.

\subsection{The trick}
We could compute $\phi$ for every point and run a linear SVM. But $\phi$ can be huge (degree-3 polynomial features of 1{,}000 inputs:
over 167 million features) or infinite. However, the dual needs only dot products, and for many $\phi$ there is a shortcut:
\[ \phi(\mathbf x)^\top\phi(\mathbf z) = x_1^2z_1^2 + 2x_1x_2z_1z_2 + x_2^2z_2^2 = (x_1z_1 + x_2z_2)^2 = (\mathbf x^\top\mathbf z)^2. \]
We never build $\phi$; we square the ordinary dot product. Any function $K(\mathbf x,\mathbf z) = \phi(\mathbf x)^\top\phi(\mathbf z)$ is a
\textbf{kernel}. Replace every $\mathbf x_i^\top\mathbf x_j$ by $K(\mathbf x_i,\mathbf x_j)$.

\begin{example}[: a kernel by hand]
$K(\mathbf x,\mathbf z) = (\mathbf x^\top\mathbf z + 1)^2$, $\mathbf x = (1,2)$, $\mathbf z = (3,-1)$: $\mathbf x^\top\mathbf z = 3 - 2 = 1$, so $K = 4$.
This equals a dot product in a 6-dimensional space of $(1, \sqrt2x_1, \sqrt2x_2, x_1^2, \sqrt2x_1x_2, x_2^2)$, computed with two
multiplications and an addition.
\end{example}

\subsection{Common kernels}
\begin{itemize}
\item \textbf{Linear}: $K = \mathbf x^\top\mathbf z$ (no mapping). Best for high-dimensional sparse data like text.
\item \textbf{Polynomial}: $K = (\gamma\mathbf x^\top\mathbf z + r)^d$: all interactions up to degree $d$.
\item \textbf{RBF (Gaussian)}: $K = \exp(-\gamma\norm{\mathbf x - \mathbf z}^2)$: a similarity, 1 for identical points, decaying with
  distance. Corresponds to an infinite-dimensional feature space. The default non-linear choice.
\item \textbf{Sigmoid}: $\tanh(\gamma\mathbf x^\top\mathbf z + r)$: not always a valid kernel.
\end{itemize}

\subsection{What makes a valid kernel}
$K$ must equal a dot product in \emph{some} space. Equivalently (Mercer's condition), for any set of points the \textbf{Gram matrix}
$[K(\mathbf x_i,\mathbf x_j)]$ must be symmetric and positive semi-definite. Sums, positive multiples and products of valid kernels are
valid, so kernels can be built from pieces (e.g. a text kernel plus a numeric kernel).

\subsection{The RBF $\gamma$}
$\gamma$ sets how far each support vector's influence reaches. Large $\gamma$: tiny bumps around each point, wiggly boundary,
overfitting (in the limit, memorisation). Small $\gamma$: smooth, nearly linear boundary, underfitting. Tune $C$ and $\gamma$ together
on a log grid (e.g. $C\in\{0.1, 1, 10, 100\}$, $\gamma\in\{0.001, 0.01, 0.1, 1\}$).

\section{Practicalities}
\begin{itemize}
\item \textbf{Scale features}: the margin is a distance; an unscaled feature in rupees dominates one in years.
\item \textbf{Cost}: kernel SVM training is between $O(n^2)$ and $O(n^3)$ and the kernel matrix needs $O(n^2)$ memory: fine for
  $10^4$, painful beyond about $10^5$ rows. Linear SVMs (liblinear, SGD with hinge loss) scale to millions. Approximate kernels
  (random Fourier features, Nystr\"om) make RBF-like models scalable.
\item \textbf{Prediction cost}: proportional to the number of support vectors (for kernel SVMs).
\item \textbf{Multiclass}: SVMs are binary; scikit-learn's \code{SVC} uses one-vs-one, \code{LinearSVC} one-vs-rest.
\item \textbf{Probabilities}: SVM outputs signed distances; Platt scaling fits a sigmoid on cross-validated scores
  (\code{probability=True}, slower).
\item \textbf{Imbalance}: \code{class\_weight} scales $C$ per class.
\item \textbf{SVR}: regression with the $\varepsilon$-insensitive tube (Chapter 4).
\item \textbf{One-class SVM}: novelty detection, a boundary around ``normal'' data.
\end{itemize}

\begin{practical}[: SVMs and kernels at InfoEdge]
Linear SVMs on TF-IDF features were the best text classifiers before deep learning and remain excellent fast baselines for
classifying job titles or resumes into the taxonomy (500 categories, millions of rows: \code{LinearSVC} trains in seconds).
The kernel idea, ``compare items by a similarity instead of explicit features'', reappears in attention (softmax of dot
products) and in the Siamese networks on the Tools slide, which learn a similarity between a job and a CV.
\end{practical}

\stuck{StatQuest: The Radial (RBF) Kernel}{\ytid{Qc5IyLW_hns}}
\section{Interview questions}

\begin{qa}{\pyq{INT: what is SVM; hinge loss, kernels, margin, support vectors} Explain SVMs in the order the interviewer wants:
margin, support vectors, hinge loss, kernels.}
\oneline{An SVM finds the separating hyperplane with the widest margin; the solution depends only on the points on or inside the
margin (support vectors); with slack it is equivalent to minimising hinge loss plus an L2 penalty; and because it only needs dot
products, kernels let it learn non-linear boundaries.}
\why
\textbf{Margin.} Among all separating hyperplanes choose the one farthest from the nearest points; scaled so the nearest satisfy
$y(\mathbf w^\top\mathbf x + b) = 1$, the margin is $2/\norm{\mathbf w}$, so we minimise $\frac12\norm{\mathbf w}^2$ subject to all points
being on the correct side with functional margin $\ge1$.
\textbf{Support vectors.} In the dual, $\mathbf w = \sum\alpha_iy_i\mathbf x_i$ with $\alpha_i > 0$ only for points on or inside the margin.
\textbf{Hinge loss.} Soft margin: $\min\frac12\norm{\mathbf w}^2 + C\sum\max(0, 1 - y_if(\mathbf x_i))$. $C$ trades margin width for
training errors.
\textbf{Kernels.} The dual and the prediction use only $\mathbf x_i^\top\mathbf x_j$; replacing with $K(\mathbf x_i,\mathbf x_j) =
\phi(\mathbf x_i)^\top\phi(\mathbf x_j)$ fits a linear separator in a high (even infinite) dimensional space without computing $\phi$.
\worked RBF SVM on a ``ring'' dataset: linear SVM 52\% accuracy; RBF with $C = 1$, $\gamma = 0.5$: 99\%.
\pushes \emph{``Why maximise the margin?''} Larger margins mean the classifier is less sensitive to small perturbations of the data
and, by statistical learning theory, generalisation bounds depend on margin rather than dimension, which is why SVMs work in very high
dimensions.
\end{qa}

\begin{qa}{\pyq{OA: hinge loss} Is hinge loss a classification or regression loss? Compute it for $(y, f) = (+1, 0.3)$,
$(+1, -2)$, $(-1, -1.5)$.}
\oneline{Classification; the values are $0.7$, $3$ and $0$.}
\why $\max(0, 1 - yf)$: $1 - 0.3 = 0.7$; $1 + 2 = 3$; $yf = 1.5$, so $0$. It is a convex upper bound on the 0--1 loss.
\end{qa}

\begin{qa}{\pyq{INT: kernels project to higher dimensions} What does the kernel trick do, exactly, and why is it a ``trick''?}
\oneline{It lets an algorithm that only uses dot products operate in a high-dimensional feature space by computing those dot products
directly with a kernel function, without ever constructing the high-dimensional vectors.}
\why ``Trick'' because the cost stays that of computing $K$ in the original space. Example: $(\mathbf x^\top\mathbf z)^2$ equals a dot
product of degree-2 features. For RBF, the implicit space is infinite-dimensional, so explicit mapping is impossible, yet
$K = e^{-\gamma\norm{\mathbf x-\mathbf z}^2}$ is cheap. Any dot-product-only algorithm can be kernelised: kernel ridge regression, kernel
PCA, kernel $k$-means, Gaussian processes.
\end{qa}

\begin{qa}{Derive why the SVM margin is $2/\norm{\mathbf w}$.}
\oneline{After scaling so the closest points satisfy $|\mathbf w^\top\mathbf x + b| = 1$, their distance to the plane is $1/\norm{\mathbf w}$,
and the street spans both sides, so its width is $2/\norm{\mathbf w}$.}
\why Distance of $\mathbf x$ to the plane is $|\mathbf w^\top\mathbf x + b|/\norm{\mathbf w}$ because $\mathbf w/\norm{\mathbf w}$ is the unit normal.
For $\mathbf w = (3,4)$: margin $2/5 = 0.4$.
\end{qa}

\begin{qa}{What does $C$ control? What happens as $C\to\infty$ and as $C\to0$?}
\oneline{$C$ is the price of margin violations: large $C$ gives a narrow margin that fits training data closely (low bias, high variance),
small $C$ a wide margin that tolerates errors (high bias, low variance); $C\to\infty$ is the hard margin, $C\to0$ ignores the data.}
\why In hinge-loss form, $C$ is like $1/\lambda$. With outliers, a large $C$ bends the boundary to accommodate them.
\end{qa}

\begin{qa}{An RBF SVM gets 100\% training accuracy and poor test accuracy. What do you change?}
\oneline{It is overfitting: decrease $\gamma$ (smoother kernel) and/or decrease $C$ (wider margin), tuned by cross-validation; also
check scaling.}
\why Large $\gamma$ makes each support vector influence only its immediate neighbourhood, so the model memorises. The number of
support vectors approaching $n$ is a symptom.
\end{qa}

\begin{qa}{Logistic regression vs linear SVM: which and when?}
\oneline{Both are linear classifiers with L2 regularisation differing only in loss; choose logistic regression for calibrated
probabilities and interpretability, linear SVM when you want a max-margin boundary that ignores well-classified points, and either for
text classification with similar accuracy.}
\why Log loss is smooth and never zero; hinge is zero beyond the margin and has a kink. SVM gives sparse dual solutions; logistic gives
probabilities. With kernels, SVM is the natural choice.
\end{qa}

\begin{qa}{Why must features be scaled for SVMs?}
\oneline{Because the margin and the RBF kernel are based on Euclidean distances, so a feature with a large numeric range dominates the
geometry regardless of its importance.}
\worked Salary in rupees (range $10^6$) and experience in years (range 30): $\norm{\mathbf x - \mathbf z}^2$ is essentially only the salary
difference.
\end{qa}

\begin{qa}{\deep Why do SVMs work well in very high dimensions with few samples (e.g. text)?}
\oneline{Because their capacity is controlled by the margin, not by the number of features, and in high dimensions data are often
nearly linearly separable with a good margin; the solution lives in the span of the support vectors, at most $n$ of them.}
\why Generalisation bounds for margin classifiers scale with $R^2/\rho^2$ (data radius over margin squared), independent of dimension.
TF-IDF vectors of documents with 100{,}000 dimensions are sparse and nearly separable; a linear SVM finds a robust boundary.
\end{qa}

\begin{qa}{\deep Can you use the kernel trick with logistic regression, $k$-means or PCA?}
\oneline{Yes: any method that can be written purely in terms of dot products can be kernelised, giving kernel logistic regression, kernel
$k$-means and kernel PCA; the catch is that without hinge loss's sparsity, all $n$ points are needed at prediction time.}
\why Kernel logistic regression: $f(\mathbf x) = \sum_i\alpha_iK(\mathbf x_i,\mathbf x)$ with log loss; every $\alpha_i\ne0$, so prediction costs
$O(n)$ and training $O(n^3)$. Kernel PCA: eigendecompose the centred kernel matrix to get non-linear components. Kernel $k$-means: distances
in feature space via $K(\mathbf x,\mathbf x) - 2K(\mathbf x,\mathbf c) + \dots$; equivalent to spectral clustering in some forms.
\end{qa}

\begin{qa}{\deep Is a neural network with a final linear layer ``just a learned kernel''?}
\oneline{In a sense yes: the hidden layers learn a feature map $\phi$, and the last layer is a linear model on $\phi(\mathbf x)$, so the
network learns its kernel $K(\mathbf x,\mathbf z) = \phi(\mathbf x)^\top\phi(\mathbf z)$ from data instead of fixing it in advance.}
\why An RBF SVM fixes the similarity; a network adapts it to the task. In the infinite-width limit, a randomly initialised network
trained by gradient descent behaves like a kernel method with the ``neural tangent kernel''. Siamese networks make this explicit: they
learn an embedding where dot products measure job--CV similarity.
\end{qa}

\begin{qa}{Compute the polynomial kernel $(\mathbf x^\top\mathbf z + 1)^2$ for $\mathbf x = (1,2)$, $\mathbf z = (3,-1)$ and the RBF kernel with
$\gamma = 0.5$ for $\mathbf x = (0,0)$, $\mathbf z = (1,1)$.}
\oneline{Polynomial: $(1+1)^2 = 4$; RBF: $e^{-0.5\cdot2} = e^{-1} = 0.368$.}
\end{qa}

\begin{qa}{What are the KKT conditions telling us about which points are support vectors?}
\oneline{Complementary slackness says $\alpha_i[y_if(\mathbf x_i) - 1 + \xi_i] = 0$, so $\alpha_i > 0$ only when a point's constraint is
tight (on the margin) or violated; points safely outside the margin have $\alpha_i = 0$.}
\why Also $(C - \alpha_i)\xi_i = 0$: points with slack ($\xi_i > 0$) have $\alpha_i = C$. So: $\alpha = 0$ outside, $0<\alpha<C$ on the margin,
$\alpha = C$ inside or wrong.
\end{qa}

\begin{qa}{What happens if you remove a non-support vector and retrain? A support vector?}
\oneline{Removing a non-support vector leaves the solution unchanged; removing a support vector can change it.}
\why This is also why leave-one-out error of an SVM is bounded by (number of support vectors)$/n$.
\end{qa}

\begin{qa}{\deep Your kernel SVM takes hours on 500{,}000 rows. What are your options?}
\oneline{Use a linear SVM if the data are high-dimensional, approximate the kernel with random Fourier features or Nystr\"om and train
a linear model on them, subsample, or switch to gradient boosting or a neural network.}
\why Kernel SVM cost grows at least quadratically. Random Fourier features: $z(\mathbf x) = \sqrt{2/D}\cos(\Omega\mathbf x + \mathbf b)$ with
$\Omega$ drawn from a Gaussian approximates the RBF kernel with $D$ features, so a linear SVM on $z$ trains in linear time.
\end{qa}

\begin{qa}{How do you get probabilities from an SVM, and why is it costly?}
\oneline{Platt scaling: fit a logistic regression mapping the SVM's decision values to probabilities, using cross-validated decision values
so the calibration data were not used to train the SVM; the internal cross-validation makes training about five times slower.}
\end{qa}

\begin{qa}{\deep When would you choose an SVM over gradient boosting today?}
\oneline{For small-to-medium data with many features and smooth or similarity-based structure (text with TF-IDF, bioinformatics,
custom kernels on strings or graphs), for strong regularised baselines, and for one-class novelty detection; for large mixed tabular data,
boosting usually wins.}
\why Kernels let you encode domain similarity (string kernels for skill sequences). SVMs have few hyperparameters and convex training
(reproducible). They lack native handling of categoricals and missing values and scale poorly.
\end{qa}

\begin{qa}{How many support vectors would you expect for a well-separated vs a noisy dataset, and what does it indicate?}
\oneline{Few for well-separated data and many for noisy or overlapping data; a very high fraction suggests overlap, too small $C$ or an
overly complex kernel (overfitting with large $\gamma$).}
\end{qa}

\begin{qa}{What is a one-class SVM and where could InfoEdge use it?}
\oneline{It learns a boundary enclosing most of the ``normal'' data (in kernel space) and flags points outside as novel; useful for
spotting unusual recruiter behaviour or anomalous job postings when labelled fraud is scarce.}
\why Parameter $\nu$ upper-bounds the fraction of training points treated as outliers. Isolation Forest is often faster on large data.
\end{qa}
